\subsection{Implications for model developers}
\label{subsec:orf-practical}
For a model generator, the contract is useful when row residuals have an application
meaning that coefficient magnitudes alone do not express. Developers should retain
the original rows and unit conversions, declare the budgets, and record whether
selection changes a proposed factor or finds incompatible requirements. Unchanged
factors still provide an explicit interpretation of accuracy; they do not demonstrate
a computational benefit from the budget.

Returned-point verification remains necessary. The exported-residual hypotheses of
the contract must be checked on the actual written representation, together with
integer rounding, bounds and original-unit residuals. A passing primal audit and the
absence of a detected bound contradiction answer narrower questions than a global
mixed-integer optimality certificate.

Deployment also requires a complete-process comparison with no external scaling and
with strong numerical controls. At the tested 100\,m\textsuperscript{3} setting,
Table~\ref{tab:base-controls-time} shows that preparation is recovered on models
001, 027 and 046, but increases total time on 005, 022 and 031. This suggests measuring
preparation and solve costs on representative workloads before enabling a policy.
It does not supply a predictive timing threshold for unseen models. Conditioning
diagnostics and the lower solver work observed here cannot replace that measurement.
